(1) 总动量为 \(m_1\boldsymbol v\)，故质心速度 \(\boldsymbol V=m_1\boldsymbol v/(m_1+m_2)\)。碰撞中没有外部平面内冲量，所以质心速度不变。

(2) 球面光滑意味着冲量仅沿 \(\boldsymbol n\)。设第一球所受冲量为 \(J\boldsymbol n\)，则 \(\boldsymbol v_1'=\boldsymbol v+J\boldsymbol n/m_1,\ \boldsymbol v_2'=-J\boldsymbol n/m_2\)。完全弹性法向恢复关系 \((\boldsymbol v_1'-\boldsymbol v_2')\cdot\boldsymbol n=-(\boldsymbol v\cdot\boldsymbol n)\)，从而
\[
J=-\frac{2m_1m_2}{m_1+m_2}(\boldsymbol v\cdot\boldsymbol n)>0.
\]
即
\(\boldsymbol v_1'=\boldsymbol v-\frac{2m_2}{m_1+m_2}(\boldsymbol v\cdot\boldsymbol n)\boldsymbol n\)，
\(\boldsymbol v_2'=\frac{2m_1}{m_1+m_2}(\boldsymbol v\cdot\boldsymbol n)\boldsymbol n\)。

(3) 用末速度展开动能差：
\[
2\Delta K=2J(\boldsymbol v\cdot\boldsymbol n)
+J^2(1/m_1+1/m_2)=0
\]
由所求 J 立即得零。冲量大小为 \(|J|\)；碰撞中的动能不因摩擦损失。

(4) 当 \(m_2\to\infty\)，\(\boldsymbol v_2'\to0\)，第一球 \(\boldsymbol v_1'=\boldsymbol v-2(\boldsymbol v\cdot\boldsymbol n)\boldsymbol n\)，法向分量反向而切向分量保持不变，恢复镜面弹性反射。
